% Independent presentation wrapper; no unknown-permission journal style/logo is copied.
\documentclass[10pt,a4paper]{article}
\usepackage[margin=22mm]{geometry}
\usepackage[english]{babel}
\usepackage{amsmath,amssymb,amsthm,euscript}
\usepackage[shortlabels,inline]{enumitem}
\usepackage[arrow,matrix,curve]{xy}
\usepackage{array,longtable,graphicx,url,multirow}
\usepackage[colorlinks]{hyperref}
\numberwithin{equation}{section}
\newtheorem{m-theorem}{Theorem}[section]
\newtheorem{m-corollary}[m-theorem]{Corollary}
\newtheorem{m-lemma}[m-theorem]{Lemma}
\newtheorem{m-proposition}[m-theorem]{Proposition}
\newtheorem{claim}{Claim}
\newtheorem*{fact}{Fact}
\newtheorem*{thm-nono}{Theorem}
\newtheorem*{prop-nono}{Proposition}
\newtheorem*{lm-nono}{Lemma}
\newtheorem*{conj-nono}{Conjecture}
\theoremstyle{definition}
\newtheorem{m-definition}[m-theorem]{Definition}
\newtheorem{m-notation}[m-theorem]{Notation}
\newtheorem{m-asmp}[m-theorem]{Assumptions}
\newtheorem*{def-nono}{Definition}
\theoremstyle{remark}
\newtheorem{m-example}[m-theorem]{Example}
\newtheorem{m-remark}[m-theorem]{Remark}
\newtheorem{m-question}[m-theorem]{Question}
\newenvironment{m-eqn}{\begin{equation}}{\end{equation}}
\setlist[enumerate,1]{label={\rm(\roman*)},ref={\rm\roman*}}
\setlist[enumerate,2]{label={\rm(\alph*)},ref={\rm\alph*}}
\newlist{aenum}{enumerate}{2}
\setlist[aenum,1]{label={\rm(\alph*)},ref={\rm\alph*}}
\setlist[aenum,2]{label={\rm(\roman*)},ref={\rm\roman*}}
\providecommand\MakeLinkTarget[2][]{}
\makeatletter\newcommand\customitem[1]{\item[#1{\MakeLinkTarget[item]{}}]\def\@currentlabel{#1}}\makeatother
\providecommand\doi[1]{\href{https://doi.org/#1}{doi:\nolinkurl{#1}}}
\providecommand\bysame{\leavevmode\hbox to 3em{\hrulefill}\,}
\let\euf\EuScript %use of package ``euscript'' required
\let\cal\mathcal
\let\mbb\mathbb
\let\mfrak\mathfrak 

\newcommand\uset[2]{{\disp\mathop{\mbox{$#2$}}_{#1}}}
\newcommand\oset[2]{{\disp\mathop{\mbox{$#2$}}^{#1}}}
\newcommand\ouset[3]{{\overset{#2}{\underset{#1}#3}\,}}

\newcommand{\Ad}{{\rm Ad}}
\newcommand{\BD}{{O_1}}
\let\dta\delta
\let\eps\epsilon
\let\ges\geqslant 
\let\les\leqslant 
\let\hra\hookrightarrow 
\let\mt\mapsto
\let\nit\noindent
\let\ovl\overline 
\let\sm\setminus 
\let\Si\Sigma 
\let\srel\stackrel 
\let\tld\tilde 
\let\unbar\underbar 
\let\vpi\varpi 
\let\vphi\varphi 
\let\disp\displaystyle 
\newcommand{\bad}{{\rm bad}}
\newcommand{\cd}{\mathop{\rm cd}\nolimits}
\newcommand{\codim}{\mathop{\rm codim}\nolimits}
\renewcommand{\dim}{\mathop{\rm dim}\nolimits}
\newcommand{\etl}{{\rm\acute{e}t}}
\newcommand{\gen}{{\rm gen}}
\newcommand{\iris}{{\rm iris}}
\newcommand{\lcit}{{\textit{loc.\,cit.}}}
\newcommand{\ocit}{{\textit{op.\,cit.}}}
\newcommand{\kk}{{\mbb C}}%{{\Bbbk}}
\newcommand{\Ker}{\mathop{\rm Ker}\nolimits}
\newcommand{\eE}{{\euf E}}
\newcommand{\eF}{{\euf F}}
\newcommand{\eG}{{\euf G}}
\newcommand{\eI}{{\euf I}}
\newcommand{\bbI}{{\mbb I}}
\newcommand{\eL}{{\euf L}}
\newcommand{\pos}{{>0}}
\newcommand{\pr}{{\rm pr}}
\newcommand{\rd}{{\rm d}}
\let\rst\vert 
\newcommand{\eN}{{\euf N}}
\newcommand{\eO}{{\euf O}}
\newcommand{\eT}{{\euf T}}
\newcommand{\cY}{{\cal Y}}
\newcommand{\bbZ}{{\mbb Z}}
\newcommand{\VV}{V}
\newcommand{\xx}{x}
\newcommand{\XX}{X}
\newcommand{\YY}{Y}
\newcommand{\tpl}{{\rm top}}
\newcommand\Hom{\mathop{\rm Hom}\nolimits}
\newcommand\Pic{\mathop{\rm Pic}\nolimits}
\newcommand{\Spec}{\mathop{\rm Spec}\nolimits}
\newcommand{\Sym}{\mathop{\rm Sym}\nolimits}
\newcommand{\res}{\mathrm{res}}
\newcommand{\proj}{\mathrm{proj}}
\newcommand{\open}{\mathrm{open}}
\newcommand{\num}{\mathrm{num}}
\newcommand{\kod}{\mathrm{kod}}
\newcommand{\rmdiv}{\mathrm{div}}
\newcommand{\amp}{\mathrm{amp}}
\DeclareMathOperator{\Grass}{Grass}
\DeclareMathOperator{\Eff}{Eff}
\DeclareMathOperator{\rank}{rank}
\DeclareMathOperator{\Sing}{Sing}

\newcommand{\lra}{\longrightarrow}
\makeatletter
\newlabel{prop:irr-fibre}{{2.10}{8}{}{}{}}
\newlabel{iti:GOX}{{\rm i}{8}{}{}{}}
\newlabel{lm:strg-gen}{{3.8}{12}{}{}{}}
\makeatother
% Technical reference only: fixed Halic arXiv2401.16196v4 native line 794 and official PDF Section 3.2.3.
\makeatletter
\newlabel{ssct:subvar-homog}{{3.2.3}{}{}{}{}}
\makeatother
\title{Mobility and irreducible general isotropy give a Chow-type G3 criterion}
\author{Mihai Halic}
\date{10.46298/epiga.2026.12974; CC BY-SA4.0}
\begin{document}
\maketitle
\noindent\textit{Editorial note.} Only fixedv4 is used; v3 ND permission is not extended or bypassed. The selected article-content portion remains CC BY-SA4.0 with attribution/change notice.\par\medskip
\noindent\textit{Editorial note.} Original mathematical spans and all current local core remain verbatim; the native686 product discrepancy is disclosed, not repaired. The existing explicit source-location note is the chosen handling; no mathematical repair is made.\par\medskip
\noindent\textit{Editorial note.} No compiler/proof checker/hash/byte comparison or author-correctness review was run.\par\medskip
\noindent\textit{Editorial note.} Source references to unselectedProp2.10(i)/Lemma3.8 retain original numbers as contextual pointers; neither substitutes for the retained core.\par\medskip
\noindent\textit{Editorial note.} Revision r002 moves only source-reference newlabel declarations (attribution lines 12–16) from the body to the preamble; the original attribution and all author mathematical text are unchanged. Failed r001 metadata is preserved locally and the root-owned failed job is untouched.\par\medskip
\noindent\textit{Editorial note.} Revision r003 supplies the exact original Section 3.2.3 label for the contextual reference at native line 414. The heading and numbering are confirmed in the fixed native source and official PDF. The section is outside the selected proof; all current core remains fully collected. No source mathematical text changes.\par\medskip
\noindent\textit{Editorial scope.} One complete fixed-version principal author claim, licensed for the selected editable packet. Scope, H2 core and bounded semantic-neighbor comparison are closed for assembly; current pre-upload identity refresh and actual compile/readability/index/export remain root-owned.\par\medskip
\noindent\textit{Allowed standard prerequisites (editorial).}\begin{itemize}
\item Classical formal geometry G2 finite-cover reduction. Badescu,Projective geometry and formal geometry(2004),Corollary9.20; cited by the author at native672 and retained in the original bibliography. External named G2-to-finite-cover theorem; not an omitted new/current local proof.
\item Stein factorization and basic finite/etale family geometry. Classical algebraic geometry; author Lemma2.9 uses Stein factorization. The entire current Lemma2.9 argument is retained. Its base-change/isotropy conclusions are not replaced by this prerequisite.
\item Speiser criterion and intersection theory. Original background native314–323, citing Hartshorne Ample Subvarieties(1970); original full bibliography; deformation-invariant intersection products. Named external equivalence used after G3 is established; local Lemma2.6 and the current ramification argument are fully retained.
\end{itemize}\medskip
\section*{Source, license and collection scope}
\noindent Mihai Halic, \emph{G3-criteria and applications}, \'Epijournal de G\'eom\'etrie Alg\'ebrique 10 (2026), Article 18, 20 pages. DOI: \href{https://doi.org/10.46298/epiga.2026.12974}{10.46298/epiga.2026.12974}. Fixed author source: \href{https://arxiv.org/abs/2401.16196v4}{arXiv:2401.16196v4}, 18 September 2026; published 21 September 2026.
\par\medskip
\noindent Copyright by the author(s). Original article and the adapted article-content portion here are distributed under \href{https://creativecommons.org/licenses/by-sa/4.0/}{CC BY-SA 4.0}. Changes: selection of Theorem 3.1 with complete author proof, necessary definitions and complete local supporting lemma proofs; independent reformatting and editorial source-location notes. Mathematical wording, formulas and proof steps in the selected author spans are unchanged. No author endorsement is implied; the license supplies no warranties and no additional reuse restrictions are imposed.
\par\medskip
\noindent One main result is collected. Its two authored cases, the cohomological-dimension conclusion and supporting lemmas are not separate counts. Main statement: native source lines 658--669; proof: 671--703. Core Lemma 2.9: 559--582. The complete current local supporting arguments are retained below, rather than substituted by prerequisite citations.
\par\medskip
\noindent\textit{Source difference retained verbatim.} Native line 686 writes $\mathcal Y'\subset\mathcal Y\times_X\mathcal Y$; the author's diagram at line 679 writes $\mathcal Y'\subset\mathcal Y\times_X X'$, and Lemma 2.9 uses the latter base-change product. Both printed forms remain as supplied; no mathematical correction or correctness certification is made here.
\par\medskip
\noindent References to the unselected source Proposition 2.10(i) and Lemma 3.8 retain their original source numbers. They appear only as contextual pointers; they are not substituted for the collected core proof. The complete fixed article remains available at \url{https://epiga.episciences.org/19406/pdf}. This material has not yet been compiled by this collector.
\setcounter{section}{1}

\section{Movable and generating subvarieties}\label{sct:gen-fam}

\subsection{Background}\label{ssct:setup}

In this article, $\XX$ stands for a complex, smooth projective variety. For a connected subvariety $\YY$ defined by the ideal $\eI_\YY\subset\eO_\XX$, let $\hat \XX_\YY$ be the formal completion of $\XX$ along $\YY$; its structure sheaf is ${\varprojlim}_m\,\eO_\XX/\eI_\YY^{m}$. The rational formal functions $\kk(\hat \XX_\YY)$ on $\XX$ along $\YY$ form a field. 
Hironaka--Matsumura~\cite{hir-mat} introduced the following generating properties: 
\begin{itemize}
\item $\YY$ is G2 in $\XX$ if $\kk(\XX)\hra \kk(\hat \XX_\YY)$ is a finite field extension. 
\item $\YY$ is G3 in $\XX$ if $\kk(\XX)\hra \kk(\hat \XX_\YY)$ is an isomorphism.
\end{itemize}
For a detailed treatment of these topics, the reader may consult the reference B\u{a}descu~\cite{bad}. The following well-known result due to Speiser relates the G2- and G3-properties.

\begin{thm-nono}[\textit{cf.}~\cite{hart-as}]
The following conditions are equivalent $($for a closed subscheme $\YY\subset\XX)$: 
\begin{enumerate}
\item $\YY$ is G3 and $\YY$ intersects all divisors of $\XX$;
\item $\YY$ is G2 and $\cd(\XX\sm \YY)\les\dim\XX-2$. 
$($Here $\cd(\cdot)$ stands for ``cohomological dimension''.$)$
\end{enumerate}
\end{thm-nono}

The author~\cite{hlc-pas} introduced a sequence of partial amplitude conditions for a subvariety, building upon work of Ottem~\cite{ottm}; the weakest was called $1^\pos$ (one-positivity). 

\begin{m-definition}[\textit{cf.}~{\cite[Proposition 1.6(iii)]{hlc-pas}}]
A locally complete intersection (lci) subvariety $\YY$ is $1^\pos$ if $\,\eN_{\YY/\XX}^\vee\!=\eI_\YY/\eI_\YY^2\;\text{is not pseudo-effective and}\;\cd(\XX\sm\YY)\les\dim\XX\!-2.$
\end{m-definition}

Recall, see~\cite[Corollary 2.6]{hlc-g2}, that the non-pseudo-effectiveness of $\eN_{\YY/\XX}^\vee$ implies that $\YY$ is $G2$ in~$\XX$. Together with Speiser's criterion, this implies that lci, $1^\pos$ subvarieties possess the G3-property. The drawback of this definition is that it requires controlling the cohomological dimension of $\XX\sm\YY$, which is not straightforward. (See~\cite{hlc+taj} for various methods to estimate it.)

To motivate the forthcoming mobility considerations, let us also mention that the author proved, see~\cite[Proposition 3.5]{hlc-g2}, that strongly movable lci subvarieties --~this is a simple geometric condition, see below~-- have non-pseudo-effective co-normal bundle. 

%%%%%%%%%%
%%%%%%%%%%

\subsection{Mobility}

For curves and divisors, the concept is classical. In higher codimensions, the notion of (very) moving subvarieties was introduced by Voisin~\cite{voisin-coniv}, in the attempt to geometrically characterize big subvarieties. A notion with similar flavour, that of generating subvarieties, appears in Chow~\cite{chow}; see also Debarre~\cite{dbar-mumford}, B\u{a}descu~\cite{bad}. The latter notion turns out to be essential for investigating the G3-property of subvarieties of homogeneous varieties. With this motivation, we consider the following. 

\begin{m-definition}\label{def:YY-gener}\leavevmode
\begin{aenum}
\item\label{d:YY-g-a} Let $S$, the parameter space, be a quasi-projective variety, and let 
\begin{equation}
\xymatrix@R=0.49em@C=3em{
&\cY\ar[dl]_-\vpi\ar[dr]^\rho &&\cY\subset S\times\XX\\ S&&\XX&}\label{eq:syx}
\end{equation}
be an irreducible, flat $S$-family of subvarieties of $\XX$; we assume that $\cY$ is normal. We consider the following varieties and schemes (they are of interest whenever non-empty):
\begin{enumerate}
\item\label{d:YY-g-a1} Let $\Si_1:=S$, $\cY_{1}=\cY$, $\vpi_1:=\vpi$, and let $\Delta_\cY\subset\cY\times\cY$ be the diagonal.
  
\item\label{d:YY-g-a2} The incidence variety $\Si_2\subset S\times S$ is the closure of the locus of the following pairs $(s,s')$: first, one has $\YY_s\cap\YY_{s'}\neq\emptyset$, $\YY_s$, $\YY_{s'}$; second, $(s,s')$ belongs to an irreducible component containing $(s,s), (s',s')$; that is,  $\YY_s\cap\YY_{s'}$ deforms (non-equidimensionally) to the ``identity intersections'' $\,\YY_s\cap\YY_s$, $\YY_{s'}\cap\YY_{s'}\,$. Below is the formal definition. 

Let $\mbb S:=\{ (s,s')\mid \YY_s=\YY_{s'} \}$; if $S$ parametrizes distinct subvarieties of $\XX$, then $\mbb S$ is the diagonal  $\Delta_S\subset S\times S$. We define $\Sigma_2$ to be the union of those irreducible components
\[
  \xymatrix{R\subset\overline{(\vpi_1,\vpi)(\cY_1\times_\XX\cY)\;\sm\;\mbb S}\subset S\times S\ar@<.35ex>[r]^-{\proj_1}\ar@<-.35ex>[r]_-{\proj_2}&S}\quad
\text{such that }
\proj_1(R)=\proj_2(R)\subset S
\](we denote the common image by $C$) and $R\supset \Delta_C$ (the diagonally embedded $C$). 
The reason for eliminating $\mbb S$ is to consider only those $\YY_s$, $s\in S$, which possess non-trivial, intersecting deformations. For movable families (see below), this definition simplifies considerably (see Remark~\ref{rmk:G}(\ref{rmk:G-2})).  

Let $\cY_2$ be defined as in the diagram below: 
\vskip-.25ex
\begin{equation}
\xymatrix@C=3.49em@R=1.79em{
\cY_2:=
\Si_2\times_{S\times S}(\cY_1\times\cY)  \ar[r]\ar@/^3ex/[rrr]|{\;\rho_2\;}\ar[d]\ar[dr]|{\;(\vpi_1,\vpi_2)\;}&
\cY_1\times\cY\ar[r]_<(.35){(\rho_1,\rho)}\ar[d]^<(.45){(\vpi_1,\vpi)}& \XX\times\XX\ar[r]_-{\proj_2\;}&\XX.
\\ 
\Si_2\ar[r]^<(.25){\iota}&S\times S&
}\label{eq:cY}
\end{equation} 
For $o\in S$, let $\{o\}\times \Sigma_o:=\iota^{-1}(\{o\}\times S)$ and $\cY_o:=\cY\rst_{\Sigma_o}$, and let $\rho_o:=\rho\rst_{\cY_o}$ --~the symbol ``$\rst$'' stands for restriction~-- be the evaluation morphism $\rho$ on those $\YY_s$ which intersect $Y_o$ non-trivially.
 
\item\label{d:YY-g-a3} Suppose $\Si_t\subset S^t$, $\cY_t=(\vpi_1,\dots,\vpi_t)^{-1}(\Si_t)\subset\cY^t$ are constructed, with morphisms 
\[
(\vpi_1,\dots,\vpi_t,\rho_1,\dots,\rho_t)\colon \cY_t\lra (S\times\dots\times S)\times (\XX\times\dots\times \XX).
\]
Here $\vpi_i$, $\rho_i$, $i=1,\dots, t$, are, respectively, the compositions of $\vpi$, $\rho$ with $\cY_t\to\cY^t\xrightarrow{\proj_i}\cY$. A~diagram similar to~\eqref{eq:cY} above allows us to define 
\begin{align*}
&\Si_{t+1}:=(\Sigma_t\times S)\times_{\Sigma_{t-1} \times S^{2}}\left(\Sigma_{t-1}\times\Sigma_2\right)\subset S^{t+1}, 
\\
&\cY_{t+1}:=\left(\vpi_1,\dots,\vpi_{t+1}\right)^{-1}\left(\Si_{t+1}\right)\subset\cY^{t+1}.
\end{align*}
There is a forgetful morphism of the $(t+1)^{\rm st}$ component $\Sigma_{t+1}\to\Sigma_{t}$, and also a ``repeating the last component'' morphism $\Sigma_t\to\Sigma_{t+1}$. We call $\Si_t$ the \emph{$t$-incidence variety of\, $S$}; roughly, it consists of sequences $(s_1,\dots,s_t)$ whose successive terms belong to $\Sigma_2$. A \emph{$\cY$-chain of length $t$} is a sequence $\YY_{s_1},\dots,\YY_{s_t}$, with $(s_1,\dots,s_t)\in\Si_t$. 
\end{enumerate}

\item \label{d:YY-g-b}
We define the following properties for the family $\cY$: 
\begin{enumerate}
\item\label{d:YY-g-b1} The family $\cY$ is \emph{movable} if $\rho\colon\cY\to \XX$ is dominant. 

\item\label{d:YY-g-b2} The family $\cY$ is \emph{strongly movable} if $(\vpi_1,\rho_2)\colon\cY_2\to S\times \XX$ is dominant and its smooth locus does intersect $(\varpi_1,\varpi_2)^{-1}(\Delta_S)$ above some open subset $S'$ of $S$. 

  In this situation, for $o\in S'$ and a general point $q\in\YY_o$, $\rd\rho_{o,q}\colon\eT_{\cY_o,(o,q)}\to\eT_{\XX,q}$ is surjective and $\rho_o$ is dominant. We say that $\YY_o$ is strongly movable. 

The family $\cY$ is \emph{strongly movable in codimension one} (codim-one, for short) if there is an open $O\subset\XX$  such that  
$O\subset\rho_2(\cY_2|_{\Sigma_2\cap (S'\times S)})$ and $\codim_\XX(\XX\sm O)\ges2$. 

\item\label{d:YY-g-b3} The family $\cY$ is \emph{generating} if $\big\{(s,x)\in S\times\XX\mid \exists t\ges 1,\;(s,x)\in(\vpi_1,\rho_t)(\cY_t)\big\}$ contains an open subset of $S\times\XX$. Equivalently, the last component of $\cY$-chains starting with $Y_o$, where $o\in S$ is general, sweeps out an open subset $O_o\subset\XX$.

\item\label{d:YY-g-b4} The family $\cY$ is \emph{chain connecting in codim-one} if the family $\cY$ is generating and the open subset above has the property $\;\codim_\XX(\XX\sm O_o)\ges2$.
\end{enumerate}

\item \label{d:YY-g-c}
We say that \emph{the smooth subvariety $Y\subset X$ is very moving} (\textit{cf.}~Voisin \cite[Definition 0.5]{voisin-coniv}) if the following condition holds:
\[
\begin{array}{l}
\forall\,[\text{generic}\;p\in\YY\;\text{and generic}\;T\subset\eT_{\XX,p}\;\text{of dimension}\;\dim\YY],\\[.5ex] 
\exists\,\text{smooth, nearby, embedded deformation $\YY'$ of $\YY$, with $p\in\YY'$, such that $\eT_{\YY',p}=T$.}
\end{array}
\]
The points $p$ constitute the \emph{very moving locus of $\YY$}.  Here we agree that the parameter space~$S$ parametrizes sufficiently many deformations $\YY'$ of $\YY$ such that the tangent spaces $\eT_{\YY', p}$, for $\YY'$ in $S_p:=\rho^{-1}(p)$, satisfy the property above. (Note that $S$ does not necessarily consist of all possible embedded deformations of $\YY$ in $\XX$. In Section~\ref{ssct:subvar-homog}, we consider only translates for a group action.) 
Voisin's condition means that 
\[
\,S_p\srel{\dta_p}{\lra} \Grass\left(\dim\YY,\eT_{\XX,p}\right),\quad \YY'\longmapsto \eT_{\YY',p} 
\]
is dominant about $\YY\in S_p$. Therefore, it remains so for $\YY'\in S$ in a neighbourhood of $\YY$; all small deformations $\YY'$ of $\YY$ are still very moving. By abuse of language, we say that \emph{the family $\cY$ is very moving} (after possibly shrinking $S$).

If $S_p$ is smooth at $\YY$, it is enough to require  $\eT_{S_p,\varpi(\YY)}\xrightarrow{\dta_{\varpi(\YY),p}}\Hom(\eT_{\YY,p},\eN_{\YY/\XX,p})$ to be surjective. Since $\eT_{S_p,\varpi(\YY)}\subset H^0(\YY,\eI_{\YY,p}\otimes\eN_{\YY/\XX})$, this map is simply the homomorphism
\[
H^0(\YY,\eI_{\YY,p}\otimes\eN_{\YY/\XX})\lra \Hom\left(\eT_{\YY,p},\eN_{\YY/\XX,p}\right),\quad \sigma\longmapsto\rd\sigma_p.
\]
Let us recall, see \cite[Remark 2.6]{voisin-coniv}, that one may assume, moreover, that $\eN_{\YY/\XX}$ is ample. This is still an open condition under deformations.

\item\label{d:YY-g-d}
The \emph{family $\cY$ is smooth at $p\in\XX$} if $\cY$ is smooth at every $(\varpi(\YY),p)\in\rho^{-1}(p)$ and $\rd\rho_{(\varpi(\YY),p)}$ is surjective: 
\[
\disp\sum_{\sigma\in\eT_{S_p,\varpi(\YY)}}\kern-2ex\rd\sigma_p\left(\eT_{\YY,p}\right)=\eN_{\YY/\XX,p}.
\]
We say that $p$ belongs to the \emph{smooth locus of\, $\cY$}. In complex-analytic language, this means that $p$ is a regular value of $\rho$. We say that $\rho$ is {\em smooth over codimension one points of\, $\XX$} --~for short, \emph{$\rho$ is smooth in codim-one}~-- if $\rho(\cY)$ contains an open subset $O$ such that $\codim_\XX(\XX\sm O)\ges2$ and $\cY$ is smooth at all $p\in O$. 

\item \label{d:YY-g-e}
  For short, we write $A^\gen$ for some open (general) subset of a variety or scheme $A$.
  \end{aenum}
\end{m-definition}


\begin{m-remark}\label{rmk:G} \leavevmode
  \begin{enumerate}
  \item\label{rmk:G-1} The formulation ``deformations of $\YY$'' refers \emph{only} to members of some family $\cY\to S$ containing $\YY$. The family is either given or suitably chosen, depending on the context. 


The flatness of $\cY\to S$ is required whenever we invoke the deformation invariance of intersection products in $H^*(X)$ with members of the family. Otherwise, the definitions make sense in general. 

\item\label{rmk:G-2} We are interested in families $\cY$ which are movable, at the very least. The differential of $\rho$ is surjective ($\rho$ is smooth) on an open locus $\cY^\gen\subset\cY$; for any point $p\in\cY^\gen$, $\cY\times_\XX\cY$ is smooth about $(p,p)$. So there is only one irreducible component of $\cY^\gen\times_\XX\cY^\gen$ containing the diagonal. Its projection to both $S$-factors is dominant; let $\mbb\YY_2$ be its closure. Thus we get a ``natural''  component $\overline{(\vpi_1,\vpi_2)(\mbb\YY_2)\sm\mbb S}\subset\overline{(\vpi_1,\vpi_2)(\mbb\YY_2)\sm\Delta_S}\subset\Sigma_2\subset S^2$. If $Y_s$, $s\in S$, are pairwise-distinct subvarieties of $\XX$, then $\mbb S=\Delta_S$.

\item\label{rmk:G-3} Note that $\Sigma_2=S\times S$ as soon as the self-intersection satisfies $[Y]^2\neq0$; any two members of the family intersect non-trivially. Thus we have $\cY_t=\cY^t$ and $\Sigma_t=S^t$ for $t\ges1$.

  The self-intersection is non-trivial in this situation: for $o\in S^\gen$, $\YY_o$ is lci, $\eN_{\YY_o/\XX}$ is ample, and  $2\dim\YY_o\!\ges\dim\XX$. The excess intersection theorem, see~\cite{fu-laz}, yields $\deg_{\eO_\XX(1)}[Y]^2\!>0$. 

\item\label{rmk:G-4} Suppose $\XX$ is homogeneous for the action of a group $G$. Chow, see~\cite{chow, dbar-mumford, bad}, introduced the notion of a generating subvariety $\YY\subset\XX$. One verifies that his definition coincides with ours for $S:=G$ and $\cY:=G\times \YY$ (see Lemma~\ref{lm:strg-gen}).
  \end{enumerate}
\end{m-remark}


\begin{m-lemma}\label{lm:mov}\leavevmode
\begin{enumerate}
\item\label{l:mov-1} For the family $\cY$, we have the following implications: 
\[
\text{movable} \;\;\Longleftarrow\;\;\text{generating}\;\;\Longleftarrow\;\;\text{strongly movable}\;\;\Longleftarrow\;\;\text{very moving.}
\]
\item\label{iti:Ygen} 
The family $\cY$ generates $\XX$ if and only if\, $(\vpi_1,\rho_t)\colon\cY_t\to S\times\XX$ is dominant for some $t\ges 2$.
\item\label{l:mov-3}  Let $o\in S$ and $p_o,q_o\in\YY_o\subset\cY$ be smooth points.  
\begin{itemize}
\item If\, $\widehat{\rd\rho}_{p_o}\colon \eT_{\cY, p_o} \to \eN_{\YY_o/\XX,p_o}$ is surjective, then $\cY$ is movable. 
\item If\, $\widehat{\rd\rho}_{p_o}\oplus\widehat{\rd\rho}_{q_o}\colon  \eT_{\cY\times_S\cY, (p_o, q_o)} \to\eN_{\YY_o/\XX,p_o}\oplus\eN_{\YY_o/\XX,q_o}$ 
is surjective, then $\cY$ is strongly movable. $($Note that $\eT_{\cY\times_S\cY, (p_o, q_o)}= \{(\xi,\xi')\mid\rd\vpi_{p_o}(\xi)-\rd\vpi_{q_o}(\xi')=0$.$)$
\end{itemize}
Here $\widehat{\rd\rho}$ is the composition of $\rd\rho$ with the projection $\eT_X\rst_{Y_o}\to\eN_{Y_o/X}$. 
\end{enumerate}
\end{m-lemma}

In general, the mobility conditions are distinct. 
Let $S$, $\YY$ be positive-dimensional, irreducible projective varieties, and consider $\cY=\XX=\YY\times S$. The morphism $\varpi$ is the natural projection, and $\rho$ is the identity. Then~$\cY$ is movable but not generating: any $\cY$-chain has length one. One checks that $\Sigma_2=\emptyset$, reflecting that the members of $\cY$ have no non-trivial intersections.

\begin{proof}
 (\ref{l:mov-1})~  The first three implications are clear. For the last, the surjectivity of $\rd\dta_{o,p}$ at general points $o\in S$, $p\in\YY_o$ implies that any $v\in\eT_{\YY_o,p}\sm\{0\}$ yields a surjective homomorphism $\rd\dta_{o,p}(v)\colon \eT_{S_p,o}\to\eN_{\YY_o/\XX,p}$. Thus, at nearby points $q\in\YY_o$, the evaluation at $q$ is surjective:  
  \[
  H^0(\eI_p\otimes\eN_{\YY_o/\XX})\supset\eT_{S_p,o}\ni\sigma\longmapsto\sigma(q)\in\eN_{\YY_o/\XX,q}.
  \]
  
(\ref{iti:Ygen})~ Each $(\vpi_1,\rho_t)(\cY_t)$ is constructible, and we have 
$S\times\XX={\bigcup}_{t\ges 1} \overline{(\vpi_1,\rho_t)(\cY_t)}$, by hypothesis.  Since $S\times\XX$ cannot be written as a countable union of proper subvarieties, for some $t$, we must have $\overline{(\vpi_1,\rho_t)(\cY_t)}=S\times\XX$. For the converse, there is nothing to prove.

(\ref{l:mov-3})~  The first statement is clear. For strong movability, the infinitesimal deformations of $(\YY_o,p_o)$ within $\cY_o$ are given by pointed subvarieties $(\YY_s,p_s)$, $p_s\in\YY_s\cap\YY_o$, and correspond to $\xi\in\eT_{\cY, p_o}$ such that $\widehat{\rd\rho}_{p_o}(\xi)=0$. The surjectivity of $\widehat{\rd\rho}_{q_o}$ yields  strong movability at the point $((o,o),(p_o,q_o))\in\cY_2$, which projects to $(o,o)\in\Delta_S$. Briefly, $(\rho,\rho)\colon \cY\times_S\cY\to\XX\times\XX$ transversely intersects  $Y_o\times\XX$ at $(p_o,q_o)$. 
\end{proof}

\begin{m-lemma}\label{lm:conn-mov}
Suppose that $\cY$ is such that the self-intersection satisfies $[Y]^2\neq0$. Then we have 
\[
\text{movable}\;\;\Longleftrightarrow\;\;\text{strongly movable.}
\] 
\end{m-lemma}

\begin{proof}
The morphism $\rho\colon \cY\to\XX$ is dominant; by Remark~\ref{rmk:G}(\ref{rmk:G-3}), we have $\cY_2=\cY^2$, $\Sigma_2=S^2$. We deduce that $(\vpi_1,\rho_2)$ is dominant.
\end{proof}

\begin{m-lemma}\label{lm:conn}
Let $\cY$ be a flat family of subvarieties, $D\subset\XX$ an irreducible divisor. The intersection product $[\YY_s]\cdot[D]$ is non-zero, so $\YY_s\cap D\neq\emptyset$ for all $s\in S$ in the following situations:
\begin{enumerate}
\item\label{l:conn-1}
$\XX$ is $\cY$-chain connected in codim-one.
\item\label{iti:Ogen} 
  $\cY$ is generating $\XX$ and $D\cap O_o\neq\emptyset$. Here $o\in S^\gen$ is such that $\YY_o\not\subset D$ and satisfies Lemma~\ref{lm:mov}\,{\rm(\ref{iti:Ygen})}: for some $t\ges 2$, $\rho_t\bigl(\cY|_{\Sigma_t\cap(\{o\}\times S^{t-1})}\bigr)$ contains the open subset $O_o$ of\, $\XX$. 
\item\label{iti:YD} 
$\cY$ is movable and $D\cap\YY_o\neq\emptyset$ for $o\in S^\gen$.
\end{enumerate}

\end{m-lemma}

\begin{proof}
(\ref{l:conn-1})~ Take $o\in S^\gen$ and suppose $[D]\cdot [\YY_o]=0$. The deformation invariance implies that, for all $s\in S$, we have $[D]\cdot[\YY_s]=0$. Then either $D\cap \YY_s=\emptyset$ or $\YY_s\subset D$; otherwise, the intersection $D\cap\YY_s$ is proper, so the intersection product does not vanish.  

 Let $O_o\subset X$ be the open subset swept out by $\cY$-chains starting with $\YY_o$. We have $\codim_\XX(\XX\sm O_o)\ges2$; thus $D\cap O_o\neq\emptyset$.  By chain connectedness, since $D\cap O_o\neq\emptyset$, there are $o=s_0,s_1,\dots,s_n\in S$ such that 
$\YY_{s_{j-1}}\cap \YY_{s_j}\neq\emptyset$ and $D\cap \YY_{s_n}\neq\emptyset$.

 The previous step implies $\YY_{s_n}\subset D$, so $D\cap \YY_{s_{n-1}} \neq\emptyset$, \textit{etc.}; inductively, we deduce $\YY_o\subset D$. Since $o\in S$ is generic, it follows that $\rho(\cY)\subset D$, contradicting that $\rho$ is dominant.

(\ref{iti:Ogen})~ 
 We repeat the argument above for a $\cY$-chain connecting the general points of $\YY_o$ and $D$. If the intersection product vanishes, it follows that $\YY_o\subset D$, and we have a contradiction. 
 
 (\ref{iti:YD})~ 
The intersection is proper, because $\cY$ is movable and $o$ is general in $S$.
\end{proof}

Chain connectedness in codimension one is important; we need conditions ensuring it. 

\begin{m-lemma}\label{lm:>0} 
Let $\cY$ be a flat, movable family of subvarieties of\, $\XX$. 
\begin{enumerate}
\item\label{iti:mori-dream} 
  Let $\XX$ be such that its effective divisors are semi-ample $($e.g.\  $\Pic(\XX)\cong\mbb Z$, or $X$ is a minimal Mori dream space$)$. Suppose $\cY$ is strongly movable. Then any member $\YY$ intersects the effective divisors of\, $\XX$, and $\cY$ is chain connecting in codim-one. 
\item \label{lm:>0-2}
If\, $\cY$ is strongly movable in codim-one, then $\cY$ is chain connecting in codim-one. 
\end{enumerate}
\end{m-lemma}

\begin{proof}
(\ref{iti:mori-dream})~ Suppose there is a $D$ with $[\YY]\cdot[D]\!=\!0$. Since $mD$ is globally generated for some $m\ges1$, $\YY$ is contained in a fibre $F$ of $\XX\srel{\phi}{\to}|mD|$; the same holds for the deformations of $\YY$. Note that $\dim |mD|\ges1$. 
Since $\YY$ is movable, after replacing $\YY$ by a deformation, we may assume that $F$ is a regular fibre of $\phi$. Then $\eN_{F/\XX}\rst_\YY$ (which is trivial, of rank at least one) is a quotient of $\eN_{\YY/\XX}$ (which is $(\dim \YY-1)$-ample, by~\cite{hlc-g2}), and we have a contradiction, so $[\YY]\cdot[D]\!\neq\!0$.

(\ref{lm:>0-2})~ Let $\cY\to\XX$, $S'\subset S$, and $O\subset \XX$ be as in Definition~\ref{def:YY-gener}(\ref{d:YY-g-b})(\ref{d:YY-g-b2}). Take $x, x'\in O$; there are $o,o' \in S'$ such that $x\in\YY_o$, $x'\in\YY_{o'}$. The hypothesis implies that there is a $p\in\YY_o$ such that $\rd\rho_p\colon \eT_{\cY_o,p}\to\eT_{\XX,p}$ is surjective, so $\rho(\cY_o)$ is open about the point $(o,p)$. In particular, the $\cY$-chains of length two starting with $\YY_o$ cover an open subset $U_o\subset\XX$. By repeating the argument, we obtain an open subset $U_{o'}$ swept out by the length two chains starting with $\YY_{o'}$. Since $U_o\cap U_{o'}\neq\emptyset$, there is a $\cY$-chain of length at most four connecting $x$ to~$x'$. 
\end{proof}

\subsection{Irreducible pointed deformation spaces and fundamental groups} \label{ssct:irr} 

\begin{m-definition}\label{def:irr-fibre}
Let $\cY\srel{(\varpi,\rho)}{\longrightarrow}S\times\XX$ be a normal, movable family of irreducible subvarieties. We say that~$\cY$ has {\em irreducible general isotropy} if  $\rho^{-1}(\xx)$ is irreducible for all points $\xx$ in a non-empty open subset $\XX^\iris_\cY\subset\rho(\cY)$. Equivalently, the geometric generic fibre of $\rho$ is connected ($\cY$ is normal). 
For $\xx\in\XX$, $S_\xx:=\rho^{-1}(\xx)=\{ s\mid \xx\in\YY_s \}$ can be viewed as a subset of $S$. The term ``isotropy'' is chosen in analogy with the case of group actions (see Proposition~\ref{prop:irr-fibre}\ref{iti:GOX}). 
\end{m-definition}


\begin{m-lemma}\label{lm:irr-fibre}\leavevmode
\begin{enumerate}
\item \label{l:if-1}
The family $\cY$ has irreducible general isotropy if and only if  $\cY\times_\XX\cY$ possesses exactly one irreducible component dominating $\XX$ $($which contains the diagonal $\Delta_\cY)$.
\item \label{l:if-2}
Suppose $\cY$ has irreducible general isotropy, and let $\XX'\srel{f}{\to}\XX$ be a finite $($surjective$)$ morphism, with $\XX'$ irreducible. Then the following properties hold true: 
\begin{itemize}
\item $\cY\times_\XX\XX'$ contains a unique irreducible component $\cY'$ which dominates $\XX'$. It contains $f^{-1}(\YY_s)$ for any $s\in S$ such that $\YY_s\cap\XX^\iris_\cY\neq\emptyset$.
\item We have $\cY\times_\XX U'\subset\cY'$, where $U'\subset\XX'$ is the \'etale locus of $f$. 
\end{itemize}
\end{enumerate}
\end{m-lemma}

\begin{proof} 
(\ref{l:if-1})~ For the necessity, note that the general fibre of $\cY\times_\XX\cY\to\XX$ is irreducible. For the sufficiency, consider the Stein factorization $\cY\to\hat\cY\to\XX$ (induced by some completion of $\cY$). Should the degree of the second map be greater than one, $\hat\cY\times_\XX\hat\cY$ possesses other components besides the diagonal, determined by pairs of distinct points in the fibres of $\hat\cY\to\XX$. Their pre-image in $\cY\times_\XX\cY$ would still dominate $\XX$.

(\ref{l:if-2})~ The general fibres of $\cY\times_\XX\XX'\to\XX'$ and $\cY\to\XX$ are isomorphic, which implies that 
\[
\cY'=\overline{{\cY'}^{\,\open}}^{\cY\times_\XX\XX'},\quad\text{where } {\cY'}^{\,\open}:=\rho^{-1}(\XX^\iris_\cY)\times_{\XX^\iris_\cY}f^{-1}(\XX^\iris_\cY).
\]
Furthermore, since $\YY_s\cap\XX^\iris_\cY\neq\emptyset$ and $f$ is finite, every component of $f^{-1}(\YY_s)$ dominates $\YY_s$, hence intersects $f^{-1}(\XX^\iris_\cY)$. 

Now proceed to the second claim. (The argument was provided by the anonymous referee.) The \'etale locus $U'\subset\XX'$ of $f$ is Zariski open, irreducible. Then $\cY\times_\XX U'\to \cY$ is \'etale, so it is an open morphism. Thus any component of $\cY\times_\XX U'$ dominates $\cY$ and hence $U'$. Therefore, $\cY\times_\XX U'$ is irreducible and is contained in $\cY'$.  
\end{proof}

\section{A Chow-type G3-criterion and applications}\label{sct:chow}

\subsection{The criterion} 
We prove that mobility conditions imply the G3-property. This result is inspired from~\cite{chow} and~\cite[Theorem 13.4]{bad}, which consider subvarieties of homogeneous varieties.


\begin{m-theorem} \label{thm:g3}
Let $\YY_o$ be a G2-subvariety of\, $\XX$ which has the following properties: 
\begin{itemize}[leftmargin=3ex]
\item It belongs to a flat, movable family $\cY={(\YY_s)}_{s\in S}$ with irreducible general isotropy. 
\item Any effective divisor of\, $\XX$ intersects $Y_s$, $s\in S^\gen$. $($Equivalently, by Lemma~\ref{lm:conn}\,{\rm(\ref{iti:YD})}, the intersection product of $\YY_o$ with any effective divisor of\, $\XX$ is numerically non-trivial.$)$
\end{itemize}
Then $\YY_o$ is G3 and $\cd(\XX\sm\YY_o)\les\dim\XX-2$ in either one of the following situations:
\begin{enumerate}
\item\label{thm:g3-1} $\XX$ is algebraically simply connected $(\pi_1^\etl(\XX)=\{1\})$.
\item\label{thm:g3-2} $[\YY_o]^2\neq0\in A^*(\XX)/\sim_{\num}$.  
\end{enumerate}
\end{m-theorem}

\begin{proof}
Since $Y_o$ is G2, there exist a normal, irreducible, projective variety $\XX'$, a G3-subvariety $\YY'_o\subset \XX'$, and a finite, surjective morphism $f\colon (\XX',\YY'_o)\to(\XX,\YY_o)$ such that $f$ is \'etale in a neighbourhood of $\YY'_o$ (\textit{cf.}~\cite[Corollary 9.20]{bad}). 
Let $D'\subset \XX'$ be the ramification divisor of $f$ and $D:=f_*(D')$; we want $D=\emptyset$, that is, $f$ is \'etale. 
Let us consider the diagram  
\[
\xymatrix@R=1.89em@C=2em{
\cY'_{\ovl{\omega}}\ar[r]\ar[d]^-{\vphi_{\ovl{\omega}}}&
\cY'_\omega\ar[r]\ar[d]^-{\vphi_\omega}
&\cY'\subset\cY\times_\XX\XX'\ar[r]^-{\rho'}\ar[d]^-\vphi&\XX'\ar[d]^-f
\\ 
\cY_{\ovl{\omega}}\ar[r]\ar[d]&\cY_\omega\ar[r]\ar[d]&\cY\ar[r]^-\rho\ar[d]^-\vpi&\XX\rlap{.}
\\ 
\ovl{\omega}:=\Spec(\ovl{\kk(S)})\ar[r]&{\omega}:=\Spec({\kk(S)})\ar[r]&S&
}
\]
Here $\cY'\subset\cY\times_\XX\cY$ is the irreducible component introduced in Lemma~\ref{lm:irr-fibre}, dominating $\XX'$. 
The ramification divisor of $\vphi\colon \cY'\to\cY$ is $D_1:=(\rho')^{-1}(D')\cap\cY'$. We are going to show that $D_1\cap \cY'_\omega=\emptyset$; it suffices to prove that $D_1\cap \cY'_{\ovl{\omega}}=\emptyset$, that is, $\vphi_{\ovl{\omega}}$ is \'etale. 

Since $f$ is \'etale about $\YY'_o$, Lemma~\ref{lm:irr-fibre} implies that $\YY'_o=\YY_o\times_\XX \YY'_o\subset\cY'$. Thus $\YY'_o$ is a component of $\cY'_o=(\vpi\vphi)^{-1}(o)$ and $\vphi$ is \'etale in its neighbourhood. But $\YY'_o$ is the specialization of a component $Z$ of the fibre $\cY'_{\ovl{\omega}}\to{\ovl{\omega}}$, so $\vphi_{\ovl{\omega}}\colon Z\to\cY_{\ovl{\omega}}$ is \'etale (because $\vphi$ is so about $\YY'_o$). The components of $\cY'_{\ovl{\omega}}$ are conjugate under the Galois group ${\rm Gal}\big(\ovl{\kk(S)}/\kk(S)\big)$, since $\cY'_\omega$ is irreducible. Thus $\vphi_{\ovl{\omega}}$ is \'etale everywhere and so is~$\vphi_\omega$. This proves that $D_1\cap\cY'_\omega=\emptyset$. 

Lemma~\ref{lm:irr-fibre} says that  $\cY'$ contains $f^{-1}(\YY_s)$ for $s\in S^\gen$. We deduce that $D_1\cap f^{-1}(\YY_s)=\emptyset$. Hence $D\cap\YY_s=\emptyset$; the subvariety $\YY_s$ does not intersect the branch divisor. This contradicts that $\YY_s\cap D\neq\emptyset$ for $s\in S^\gen$ (or that $[\YY_o]\cdot[D]\neq0$ ) unless $D=\emptyset$, so $f$ is \'etale.

(\ref{thm:g3-1})~ If $\XX$ is algebraically simply connected, then $f$ is an isomorphism, so $\YY_o$ is G3 in $\XX$. 

(\ref{thm:g3-2})~ Now suppose $d:=\deg(f)\ges2$, and let $\YY_o',\YY_1',\dots,\YY_{d-1}'$ be the pre-images of $\YY_o$ by $f$; they are deformation equivalent.  Then $[\YY_o]^2$ can be computed on the covering space $\XX'$: 
$$
d\cdot[\YY_o]^2=[\YY_o']\cdot [f^{-1}(\YY_o)]
=[\YY_o']\cdot([\YY_o']+[\YY_1']+\dots+[\YY_{d-1}'])
=d\cdot[\YY_o']\cdot[\YY_1']=0.
$$
\nit This contradicts the hypothesis, so $d=1$ and $f$ is \'etale.
In both situations, $Y_o$ is G3 and intersects all effective divisors, so Speiser's criterion implies that $\cd(\XX\sm\YY_o)\les\dim\XX-2$. 
\end{proof}

\newcommand{\etalchar}[1]{$^{#1}$}
\begin{thebibliography}{SVdV86+++}

\bibitem[Bad04]{bad}
L.~B\u{a}descu, \emph{Projective geometry and formal geometry}, IMPAN Monogr.\ Mat.\ (N.S.), vol.~65, Birkh\"auser Verlag, Basel, 2004, \doi{10.1007/978-3-0348-7936-1}.

\bibitem[Bad09]{bad-deb}
\bysame,  \emph{On a Connectedness Theorem of Debarre}, in: \emph{Interactions of Classical and Numerical Algebraic Geometry: A Conference in Honor of Andrew~Sommese}, D.~Bates \emph{et al.} (eds), pp.~11--20, Contemp.\ Math., vol.~496, Amer.\ Math.\ Soc., Providence, RI, 2009, \doi{10.1090/conm/496/09716}.

\bibitem[Bor91]{bor}
A.~Borel, \emph{Linear Algebraic Groups}, 2nd ed., Grad.\  Texts in Math., vol.~126, Springer-Verlag, New York, 1991, \doi{10.1007/978-1-4612-0941-6}.

\bibitem[BDP\etalchar{+}13]{bdpp}
S.~Boucksom, J.-P.~Demailly, M.~P\u{a}un, and T.~Peternell, \emph{The pseudo-effective cone of a compact K\"ahler manifold and varieties of negative Kodaira dimension}, J.~Algebraic Geom.\ \textbf{22} (2013), no.~2, 201--248, \doi{10.1090/S1056-3911-2012-00574-8}.

\bibitem[Cho69]{chow}
W.\,L.~Chow, \emph{On meromorphic maps of algebraic varieties}, Ann.\ of Math.~(2) \textbf{89} (1969), no.~2, 391--403, \doi{10.2307/1970675}.

\bibitem[Deb96]{dbar-mumford}
O.~Debarre,  \emph{Sur un th\'eor\`eme de connexit\'e de Mumford pour les espaces homog\`enes}, Manuscripta Math.\ \textbf{89} (1996), no.~4, 407--425, \doi{10.1007/BF02567526}.

\bibitem[DEL\etalchar{+}11]{delv}
 O.~Debarre, L.~Ein, R.~Lazarsfeld, and C.~Voisin, \emph{Pseudoeffective and nef classes on abelian varieties}, Compos.\ Math.\ \textbf{147} (2011), no.~6, 1793--1818, \doi{10.1112/S0010437X11005227}.

\bibitem[Fal81a]{falt-krit}
G.~Faltings, \emph{Ein Kriterium f\"ur vollst\"andige Durchschnitte}, Invent.\ Math.\ \textbf{62} (1981),  no.~3, 393--401, \doi{10.1007/BF01394251}.

\bibitem[Fal81b]{falt-homog}
  \bysame,  \emph{Formale Geometrie und homogene R\"aume}, Invent.\ Math.\ \textbf{64} (1981), no.~1, 123--165, \doi{10.1007/BF01393937}.

\bibitem[FL17]{ful+leh}
M.~Fulger and B.~Lehmann,  \emph{Positive cones of dual cycle classes}, Alg.\ Geom.\ \textbf{4} (2017), no.~1, 1--28, \doi{10.14231/AG-2017-001}.

\bibitem[FL82]{fu-laz}
W.~Fulton and R.~Lazarsfeld, \emph{Positivity and excess intersection} in: \emph{Enumerative Geometry and Classical Algebraic Geometry}, Le Barz \textit{et al.} (eds), pp.~97--105, Progr.\ Math., vol.~24, Birkh\"auser, Boston, MA, 1982, \doi{10.1007/978-1-4684-6726-0\_5}.

\bibitem[Gol82]{gold}
N.~Goldstein, \emph{Ampleness and connectedness in complex $G/P$}, Trans.\ Amer.\ Math.\ Soc.\ \textbf{274} (1982), no.~1, 361--373, \doi{10.2307/1999515}.

\bibitem[Hal19]{hlc-g2}
M.~Halic, \emph{Subvarieties with partially ample normal bundle}, Math.~Z.\ \textbf{293} (2019), no.~1-2, 371--382, \doi{10.1007/s00209-018-2175-1}.

\bibitem[Hal20]{hlc-pas}
\bysame,  \emph{Partially ample subvarieties of projective varieties}, Eur.~J.\ Math.\ \textbf{6} (2020), no.~2, 400--419, \doi{10.1007/s40879-019-00325-8}.

\bibitem[HT17]{hlc+taj}
M.~Halic and R.~Tajarod, \emph{About the cohomological dimension of certain stratified varieties}, Proc.\ Amer.~Math.\ Soc.\ \textbf{145} (2017), no.~12, 5157--5167, \doi{10.1090/proc/13691}.

\bibitem[Har68]{hart-cd}
R.~Hartshorne,  \emph{Cohomological dimension of algebraic varieties}, Ann.\ of Math.~(2) \textbf{88} (1968), no.~3, 403--450, \doi{10.2307/1970720}.

\bibitem[Har70]{hart-as}
  \bysame, \emph{Ample Subvarieties of Algebraic Varieties} (notes written in collaboration with C.~Musili), Lecture Notes in Math., vol.~156, Springer-Verlag, Berlin-New York, 1970.

\bibitem[Har74]{hart-conj}
  \bysame, \emph{Varieties of small codimension in projective space}, Bull.\ Am.\ Math.\ Soc.\ \textbf{80} (1974), no.~6, 1017--1032, \doi{10.1090/S0002-9904-1974-13612-8}.

\bibitem[HM68]{hir-mat}
H.~Hironaka and H.~Matsumura, \emph{Formal functions and formal embeddings}, J.~Math.\ Soc.\ Japan \textbf{20} (1968), no.~1-2, 52--82,\doi{10.2969/jmsj/02010052}.

\bibitem[HK00]{hu-keel}
Y.~Hu and S.~Keel,  \emph{Mori dream spaces and GIT}, Michigan Math.~J.\ \textbf{48} (2000), no.~1, 331--348, \doi{10.1307/mmj/1030132722}.

\bibitem[May90]{may}
J.\,P.~May,  \emph{Weak equivalences and quasifibrations} in: \emph{Groups of Self-Equivalences and Related Topics}, R.~Piccinini (ed.), pp.~91--101, Lecture Notes in Math., vol.~1425, Springer-Verlag, Berlin, 1990, \doi{10.1007/BFb0083834}.

\bibitem[NR98]{nap-ram}
T.~Napier and M.~Ramachandran,  \emph{The $L^2\;\bar\partial$-method, weak Lefschetz theorems, and the topology of K\"ahler manifolds}, J.~Amer.\ Math.\ Soc.\ \textbf{11} (1998), no.~2, 375--396, \doi{10.1090/S0894-0347-98-00257-4}.

\bibitem[Ott12]{ottm}
J.\,C.~Ottem, \emph{Ample subvarieties and $q$-ample divisors}, Adv.\ Math.\ \textbf{229} (2012), no.~5, 2868--2887, \doi{10.1016/j.aim.2012.02.001}.

\bibitem[Pop07]{pop}
V.~Popov, \emph{Generically multiple transitive algebraic group actions} in: \emph{Algebraic Groups and Homogeneous Spaces}, V.~Mehta (ed.), pp.~481--523, Tata Inst.\ Fund.\ Res.\ Stud.\ Math.\ \textbf{19},  Tata Inst.\ Fund.\ Res., Mumbai, 2007.

\bibitem[Sen15]{sena}
C.~Senapathi, \emph{Theorems of Barth-Lefschetz type and Morse theory on the space of paths in homogeneous spaces}, Geom.\ Dedicata \textbf{178} (2015), no.~1, 195--217, \doi{10.1007/s10711-015-0053-0}.

\bibitem[SVdV86]{soms-vdv}
A.\,J.~Sommese and A.~Van de Ven, \emph{Homotopy groups of pullbacks of varieties}, Nagoya Math.~J.\ \textbf{102} (1986), 79--90, \doi{10.1017/S002776300000043X}.

\bibitem[Ver76]{ver}
J.-L.~Verdier, \emph{Stratifications de Whitney et th\'eo\`eme de Bertini-Sard}, Invent.\ Math.\ \textbf{36} (1976), no.~1, 295--312, \doi{10.1007/BF01390015}.

\bibitem[Voi10]{voisin-coniv}
C.~Voisin, \emph{Coniveau $2$ complete intersections and effective cones}, Geom.\  Funct.\ Anal.\ \textbf{19} (2010), no.~5, 1494--1513, \doi{10.1007/s00039-009-0036-0}.

\end{thebibliography}

\end{document}
